% Luc Destombe --- licence CC BY-NC-SA 4.0
% https://creativecommons.org/licenses/by-nc-sa/4.0/deed.fr
% Fichier autonome : compiler avec pdflatex, deux fois (signets, nombre de pages).
\documentclass[a4paper,10pt]{article}
\makeatletter
\newif\iflycee@niveaux
\lycee@niveauxtrue
\newif\iflycee@palatino
\lycee@palatinotrue

\RequirePackage[utf8]{inputenc}
\RequirePackage[T1]{fontenc}
\RequirePackage[french]{babel}
\RequirePackage{etoolbox}

\newcommand{\ShorthandsOff}{\@ifundefined{shorthandoff}{}{\shorthandoff{;:!?}}}
\newcommand{\ShorthandsOn}{\@ifundefined{shorthandon}{}{\shorthandon{;:!?}}}
\ShorthandsOff
\AtBeginDocument{\ShorthandsOn}
\AtBeginEnvironment{tikzpicture}{\ShorthandsOff}

\RequirePackage{geometry}
\RequirePackage{fancyhdr}
\RequirePackage{lastpage}
\RequirePackage{multicol}
\RequirePackage{array}
\RequirePackage{enumitem}
\RequirePackage{xcolor}
\RequirePackage{needspace}

\RequirePackage{amsmath,amssymb}
\RequirePackage{mathpazo}
\DeclareMathSizes{10.95}{10.95}{8}{5.5}
\begingroup\catcode`\_=\active \catcode`\^=\active
\gdef_#1{\sb{\scriptscriptstyle #1}}
\gdef^#1{\sp{\scriptscriptstyle\lycee@moinsepais #1}}
\endgroup
\newcommand\lycee@moins{\mathbin{%
  \setbox\z@\hbox{$\m@th\scriptscriptstyle\lycee@moinsorig$}%
  \dimen@=.55\wd\z@ \advance\dimen@-.5pt
  \hbox to.75\wd\z@{\hss$\m@th\scriptscriptstyle\vcenter{\hbox{%
    \tikz\draw[line width=.5pt,line cap=round](0,0)--(\the\dimen@,0);}}$\hss}}}
\begingroup\lccode`\~=`\- \lowercase{\endgroup\let~}\lycee@moins
\newcommand\lycee@moinsepais{\mathcode`\-="8000 }
\AtBeginDocument{\mathchardef\lycee@moinsorig=\mathcode`\-\relax}
\AtBeginDocument{\catcode`\_=12 \mathcode`\_="8000
  \catcode`\^=12 \mathcode`\^="8000 }
\newcommand\lycee@exposants{%
  \fontdimen13\textfont2=.48\fontdimen6\textfont2
  \fontdimen14\textfont2=.48\fontdimen6\textfont2
  \fontdimen15\textfont2=.48\fontdimen6\textfont2
  \fontdimen13\scriptfont2=.48\fontdimen6\scriptfont2
  \fontdimen14\scriptfont2=.48\fontdimen6\scriptfont2
  \fontdimen15\scriptfont2=.48\fontdimen6\scriptfont2}
\every@math@size\expandafter{\the\every@math@size\lycee@exposants}
\newlength{\retraitcalculs}
\setlength{\retraitcalculs}{2em}

\RequirePackage{tikz}
\usetikzlibrary{arrows.meta,calc,babel,shapes.misc}
\RequirePackage[most]{tcolorbox}
\tcbuselibrary{skins,breakable}

\definecolor{coulDef}{HTML}{1F4E79}
\definecolor{coulProp}{HTML}{7030A0}
\definecolor{coulRem}{HTML}{C55A11}
\definecolor{coulEx}{HTML}{2E7D32}
\definecolor{coulExo}{HTML}{B03A2E}
\definecolor{coulPart}{HTML}{1F4E79}
\definecolor{coulMeth}{HTML}{1B6E4F}
\definecolor{coulCourbe}{HTML}{0B5ED7}
\definecolor{coulCourbeB}{HTML}{C00000}
\definecolor{coulCourbeC}{HTML}{2E7D32}
\definecolor{coulGrille}{HTML}{8C8C8C}
\definecolor{coulRep}{HTML}{1B6E4F}
\definecolor{coulBar}{HTML}{2E7D32}

\newcommand{\R}{\mathbb{R}}
\newcommand{\Rs}{\mathbb{R}^{*}}
\renewcommand{\le}{\leqslant}
\renewcommand{\ge}{\geqslant}
\renewcommand{\approx}{\simeq}
\newcommand{\vect}[1]{\overrightarrow{#1}}
\newcommand{\norme}[1]{\left\lVert #1 \right\rVert}
\newcommand{\intervalle}[2]{\left[#1\,;#2\right]}
\RequirePackage{cancel}
\renewcommand{\CancelColor}{\color{coulExo}}
\newcommand{\fois}[1]{\times{\color{coulExo}#1}}
\newcommand{\barre}[1]{\cancel{#1}}

\tikzset{grille/.style={coulGrille,line width=0.4pt}}
\newcommand{\quadrillage}[4]{\draw[grille,step=1] (#1,#2) grid (#3,#4);}
\tikzset{
  croix/.style={cross out,draw,line width=0.6pt,inner sep=0pt,minimum size=4pt},
  croixdroite/.style={croix,rotate=45,line width=0.8pt},
}
\newcommand{\pt}[3]{\path (#1) node[croix]{} node[#2,font=\small,inner sep=2.5pt]{$#3$};}

\newcommand{\lycee@repere}[4]{%
  \draw[grille,step=1] (#1,#2) grid (#3,#4);
  \draw[-{Stealth[length=2mm]},thick] (#1,0)--({#3+0.4},0) node[below left=-1pt]{$x$};
  \draw[-{Stealth[length=2mm]},thick] (0,#2)--(0,{#4+0.4}) node[below left=-1pt]{$y$};
  \node[below left,font=\scriptsize,inner sep=1.5pt] at (0,0) {$0$};}
\newcommand{\lycee@gradx}[1]{\foreach \k/\l in {#1}{%
  \draw (\k,0.07)--(\k,-0.07) node[below,font=\scriptsize,inner sep=1.5pt]{$\l$};}}
\newcommand{\lycee@grady}[1]{\foreach \k/\l in {#1}{%
  \draw (0.07,\k)--(-0.07,\k) node[left,font=\scriptsize,inner sep=1.5pt]{$\l$};}}
\newcommand{\lycee@intff}[2]{\left[#1\,;#2\right]}
\newcommand{\lycee@intfo}[2]{\left[#1\,;#2\right[}
\newcommand{\lycee@intof}[2]{\left]#1\,;#2\right]}
\newcommand{\lycee@intoo}[2]{\left]#1\,;#2\right[}
\newcommand{\lycee@ens}[1]{\left\{#1\right\}}
\AtBeginDocument{%
  \providecommand{\repere}{\lycee@repere}%
  \providecommand{\gradx}{\lycee@gradx}%
  \providecommand{\grady}{\lycee@grady}%
  \providecommand{\Cf}{\mathcal{C}\sb{\scriptscriptstyle f}}%
  \providecommand{\Cg}{\mathcal{C}\sb{\scriptscriptstyle g}}%
  \providecommand{\intff}{\lycee@intff}%
  \providecommand{\intfo}{\lycee@intfo}%
  \providecommand{\intof}{\lycee@intof}%
  \providecommand{\intoo}{\lycee@intoo}%
  \providecommand{\ens}{\lycee@ens}}

\newlist{questions}{enumerate}{3}
\setlist[questions,1]{label=\arabic*\textdegree),leftmargin=*,itemsep=2pt,topsep=3pt}
\setlist[questions,2]{label=\alph*),leftmargin=*,itemsep=1pt,topsep=2pt}
\setlist[questions,3]{label=\roman*),leftmargin=*,itemsep=1pt,topsep=1pt}
\setlist[itemize]{label=\textbullet,leftmargin=*,itemsep=2pt,topsep=3pt}

\newcommand{\besoinplace}[1]{\par\penalty\@M\begingroup
  \dimen@=#1\relax\dimen@ii=\pagegoal\advance\dimen@ii-\pagetotal
  \ifdim\dimen@>\dimen@ii\ifdim\dimen@ii>\z@\vfil\fi\break\fi\endgroup}

\newcommand{\lycee@pied}{\thepage/\pageref{LastPage}}
\newcommand{\entete}[2]{%
  \pagestyle{fancy}\fancyhf{}%
  \fancyhead[L]{\footnotesize\color{coulPart}\textbf{#1}}%
  \fancyhead[R]{\footnotesize\color{coulPart}\textbf{#2}}%
  \fancyfoot[C]{\footnotesize\lycee@pied}%
  \renewcommand{\headrulewidth}{0.4pt}%
  \renewcommand{\headrule}{\hbox to\headwidth{%
    \color{coulPart!40}\leaders\hrule height \headrulewidth\hfill}}}

\RequirePackage{ccicons}
\newcommand{\lycee@licence}{{\scriptsize\color{gray}%
  \copyright\ Luc Destombe --- \ccbyncsa\ CC BY-NC-SA 4.0}}
\newif\iflycee@licence \lycee@licencetrue
\newcommand{\sanslicence}{\lycee@licencefalse}
\AtBeginDocument{\iflycee@licence\AddToHookNext{shipout/foreground}{%
  \put(0,0){\rlap{\kern\dimexpr1in+\hoffset+\oddsidemargin+\textwidth\relax
    \llap{\raisebox{-\dimexpr1in+\voffset+\topmargin+\headheight+\headsep
      +\textheight+\footskip\relax}{\lycee@licence}}}}}\fi}

\newcommand{\formulesserrees}{%
  \setlength{\abovedisplayskip}{3pt}\setlength{\belowdisplayskip}{3pt}%
  \setlength{\abovedisplayshortskip}{2pt}\setlength{\belowdisplayshortskip}{3pt}}

\setlength{\parindent}{0pt}

\RequirePackage[implicit=false,hidelinks,pdfencoding=auto,psdextra,bookmarksopen,
  bookmarksopenlevel=1,pdfcreator={},pdfproducer={}]{hyperref}
\RequirePackage{bookmark}
\hypersetup{pdfauthor={Luc Destombe},pdfkeywords={CC BY-NC-SA 4.0}}
\newcounter{lycee@signet}
\newcommand{\signet}[3][]{\stepcounter{lycee@signet}%
  \edef\lycee@dest{\if\relax\detokenize{#1}\relax
    signet.\arabic{lycee@signet}\else#1\fi}%
  \hypertarget{\lycee@dest}{}\bookmark[level=#2,dest=\lycee@dest]{#3}}
\pdfstringdefDisableCommands{%
  \def\R{R}\def\vect#1{#1}\def\Cf{Cf}\def\Cg{Cg}\def\fois#1{×#1}%
  \def\barre#1{#1}\def\intervalle#1#2{[#1;#2]}\def\intff#1#2{[#1;#2]}%
  \def\textsuperscript#1{#1}\def\dfrac#1#2{#1/#2}\def\frac#1#2{#1/#2}%
  \def\vec#1{#1⃗}}%
\RequirePackage[official]{eurosym}

\geometry{top=0.8cm,bottom=1.3cm,left=1cm,right=1cm,footskip=0.6cm}

\pagestyle{fancy}
\fancyhf{}
\renewcommand{\headrulewidth}{0pt}
\fancyfoot[C]{\footnotesize Page \thepage{} / \pageref{LastPage}}

\newcommand{\titrefiche}[2]{%
  \begin{center}
    \Large\textbf{#1}\\[2pt]
    \large\textbf{#2}\\[2pt]
  \end{center}
  \vspace{0.10cm}}

\setlength{\columnseprule}{0.4pt}
\setlength{\columnsep}{15pt}
\renewcommand{\columnseprulecolor}{\color{black!45}}
\setlength{\multicolsep}{2pt}

\newcounter{partiefiche}
\renewcommand{\thepartiefiche}{\Roman{partiefiche}}
\newif\ifapresbandeau
\newcommand{\lycee@bandeau}[1]{%
  \par\penalty-600
  \vspace{0.08cm}%
  \noindent\colorbox{coulPart}{%
    \parbox{\dimexpr\linewidth-2\fboxsep\relax}{%
      \color{white}\bfseries\raggedright\strut#1\strut}}%
  \nopagebreak\par\nopagebreak
  \vspace{0.06cm}\nopagebreak
  \apresbandeautrue}
\newcommand{\partiefiche}[1]{%
  \refstepcounter{partiefiche}%
  \lycee@bandeau{\signet{1}{\thepartiefiche{} --- #1}\thepartiefiche{} --- #1}}
\newcommand{\partiefichesansnum}[1]{\lycee@bandeau{\signet{1}{#1}#1}}

\newcounter{exo}
\newlength{\espaceexo}\setlength{\espaceexo}{7pt}
\newlength{\espacetitre}\setlength{\espacetitre}{2pt}
\newcommand{\exo}[1][\relax]{%
  \par
  \ifapresbandeau\apresbandeaufalse\else\penalty-150\addvspace{\espaceexo}\fi
  \refstepcounter{exo}%
  \noindent\signet[exercice-\arabic{exo}]{2}{Exercice \arabic{exo}}%
  \textbf{\color{coulExo}\underline{Exercice \theexo}}%
  \ifx\relax#1\relax\else\textbf{ : #1}\fi
  \par\nopagebreak\vspace{\espacetitre}\nopagebreak
  \ignorespaces}

\setlist[enumerate]{nosep,leftmargin=*,itemsep=2pt,topsep=2pt}
\setlist[itemize]{nosep,leftmargin=*,topsep=2pt}
\newcommand{\choix}[3]{\par\hspace*{1em}%
  \makebox[0.3\linewidth][l]{a) #1}\makebox[0.3\linewidth][l]{b) #2}c) #3}
\newcommand{\deux}[2]{\par\hspace*{0.6em}%
  \makebox[0.48\linewidth][l]{#1}#2}
\newcommand{\trois}[3]{\par\hspace*{0.6em}%
  \makebox[0.32\linewidth][l]{#1}\makebox[0.32\linewidth][l]{#2}#3}

  \renewcommand{\exo}[1][\relax]{%
    \ifapresbandeau\apresbandeaufalse\else\par\penalty-150\fi
    \refstepcounter{exo}%
    \vspace{0.05cm}%
    \noindent\signet[exercice-\arabic{exo}]{2}{Exercice \arabic{exo}}%
    \textbf{\color{coulExo}\underline{Exercice \theexo}}%
    \ifx\relax#1\relax\else\textbf{ : #1}\fi%
    \par\nopagebreak
    \ignorespaces}
  \newcommand{\rappel}[1]{%
    \par\vspace{0.06cm}\nopagebreak
    \noindent\fcolorbox{black!35}{black!5}{%
      \parbox{\dimexpr\linewidth-2\fboxsep-2\fboxrule\relax}{%
        \small\textbf{Rappel.} #1}}%
    \par\vspace{0.06cm}\nopagebreak}
  \setlist[enumerate]{nosep,leftmargin=*}
  \setlist[itemize]{nosep,leftmargin=*}
  \setlength{\multicolsep}{3pt plus 1pt minus 1pt}
  \setlength{\premulticols}{10pt}
  \setlength{\postmulticols}{10pt}
  \newlist{expr}{enumerate}{1}
  \setlist[expr]{label={\Alph*\ $=$},nosep,leftmargin=*,itemsep=0pt}

\renewenvironment{center}{\par\addvspace{3pt}\centering}{\par\addvspace{3pt}}
\formulesserrees
\AtBeginDocument{\formulesserrees}
\clubpenalty=10000
\widowpenalty=10000
\displaywidowpenalty=10000
\brokenpenalty=10000
\setlength{\parskip}{0pt}

\RequirePackage{ragged2e}
\setlength{\RaggedRightRightskip}{0pt plus 1fil}
\AtBeginDocument{\RaggedRight\binoppenalty=10000 \relpenalty=10000 }
\g@addto@macro\@parboxrestore{\RaggedRight}
\makeatother
\geometry{top=0.8cm,bottom=1.3cm,left=1cm,right=1cm,footskip=0.7cm,headheight=12pt}
\setlength{\espaceexo}{5pt}

\tikzset{evt/.style={inner sep=2pt},
  proba/.style={font=\footnotesize,fill=white,inner sep=1pt},
  nom/.style={font=\footnotesize,inner sep=1.5pt,sloped},
  arbA/.style={},arbnA/.style={},arbAB/.style={},arbAnB/.style={},arbnAB/.style={},
  arbnAnB/.style={},chemin/.style={coulExo,line width=1.4pt}}
\newcommand{\arbcontraires}[2]{\def\arbnA{#1}\def\arbnB{#2}}
\newcommand{\arbfeuilles}[4]{\def\arbfAB{#1}\def\arbfAnB{#2}\def\arbfnAB{#3}\def\arbfnAnB{#4}}
\newcommand{\arbsansfeuilles}{\def\arbpasfeuilles{}}
\newcommand{\arbvaleurs}{\def\arbseules{}\arbsansfeuilles}
\newcommand{\arblab}[2]{\ifx\arbseules\undefined
  $#1\if\relax\detokenize{#2}\relax\else=#2\fi$\else$#2$\fi}
\newcommand{\arbfeuille}[2]{%
  \node[right,font=\footnotesize] at (#1.east) {$#2\ifx\arbtmp\empty\else=\arbtmp\fi$};}
\newcommand{\arbre}[8]{%
  \ifx\arbnA\undefined\def\arbNA{\overline{#1}}\else\let\arbNA\arbnA\fi
  \ifx\arbnB\undefined\def\arbNB{\overline{#2}}\else\let\arbNB\arbnB\fi
  \node (O) at (0,0) {};
  \node[evt] (A) at (3.4,1.2) {$#1$};
  \node[evt] (nA) at (3.4,-1.2) {$\arbNA$};
  \node[evt] (AB) at (7.2,1.9) {$#2$};
  \node[evt] (AnB) at (7.2,0.5) {$\arbNB$};
  \node[evt] (nAB) at (7.2,-0.5) {$#2$};
  \node[evt] (nAnB) at (7.2,-1.9) {$\arbNB$};
  \draw[arbA] (O.center) -- node[nom,above]{\arblab{P(#1)}{#3}} (A);
  \draw[arbnA] (O.center) -- node[nom,below]{\arblab{P(\arbNA)}{#4}} (nA);
  \draw[arbAB] (A) -- node[nom,above]{\arblab{P_{#1}(#2)}{#5}} (AB);
  \draw[arbAnB] (A) -- node[nom,below]{\arblab{P_{#1}(\arbNB)}{#6}} (AnB);
  \draw[arbnAB] (nA) -- node[nom,above]{\arblab{P_{\arbNA}(#2)}{#7}} (nAB);
  \draw[arbnAnB] (nA) -- node[nom,below]{\arblab{P_{\arbNA}(\arbNB)}{#8}} (nAnB);
  \ifx\arbpasfeuilles\undefined
    \ifx\arbfAB\undefined\arbfeuilles{}{}{}{}\fi
    \let\arbtmp\arbfAB\arbfeuille{AB}{P(#1\cap #2)}
    \let\arbtmp\arbfAnB\arbfeuille{AnB}{P(#1\cap\arbNB)}
    \let\arbtmp\arbfnAB\arbfeuille{nAB}{P(\arbNA\cap #2)}
    \let\arbtmp\arbfnAnB\arbfeuille{nAnB}{P(\arbNA\cap\arbNB)}
  \fi}
\newcommand{\feuille}[2]{\node[right,font=\footnotesize] at (#1.east) {#2};}
\newcommand{\arbremachines}[3]{%
  \node (O) at (0,0) {};
  \foreach \p/\q/\r [count=\i] in {#3}{%
    \pgfmathsetmacro{\y}{4.6-2.3*\i}
    \node[evt] (M\i) at (3.4,\y) {$#1_{\i}$};
    \ifnum\i=3 \tikzset{arbpos/.style={below}}\else\tikzset{arbpos/.style={above}}\fi
    \draw (O.center) -- node[nom,arbpos]{$P(#1_{\i})=\p$} (M\i);
    \node[evt] (D\i) at (7.2,\y+0.6) {$#2$};
    \node[evt] (nD\i) at (7.2,\y-0.6) {$\overline{#2}$};
    \draw (M\i) -- node[nom,above]{$P_{#1_{\i}}(#2)=\q$} (D\i);
    \draw (M\i) -- node[nom,below]{$P_{#1_{\i}}(\overline{#2})={\color{coulExo}\r}$} (nD\i);
    \node[right,font=\footnotesize] at (D\i.east) {$P(#1_{\i}\cap #2)$};
    \node[right,font=\footnotesize] at (nD\i.east) {$P(#1_{\i}\cap\overline{#2})$};}}
\newcommand{\arbretrois}[2]{%
  \node (O) at (0,0) {};
  \foreach \a/\ya/\na/\pa in {1/1.8/S/#1,0/-1.8/\overline{S}/#2}{%
    \node[evt] (N\a) at (2.2,\ya) {$\na$};
    \draw (O.center) -- node[proba]{$\pa$} (N\a);
    \foreach \b/\db/\nb/\pb in {1/0.9/S/#1,0/-0.9/\overline{S}/#2}{%
      \pgfmathsetmacro{\yb}{\ya+\db}
      \node[evt] (N\a\b) at (4.4,\yb) {$\nb$};
      \draw (N\a) -- node[proba]{$\pb$} (N\a\b);
      \foreach \c/\dc/\nc/\pc in {1/0.45/S/#1,0/-0.45/\overline{S}/#2}{%
        \pgfmathsetmacro{\yc}{\yb+\dc}
        \node[evt] (N\a\b\c) at (7,\yc) {$\nc$};
        \draw (N\a\b) -- node[proba]{$\pc$} (N\a\b\c);}}}}

\begin{document}
\titrefiche{1\textsuperscript{re} spécialité --- Probabilités conditionnelles}{Fiche d'exercices du chapitre}

\begin{multicols}{2}

\partiefiche{Vocabulaire des probabilités (rappels de seconde)}

\exo[Décrire des événements]
Un jeu en ligne tire au hasard un nombre entier de $1$ à $24$. On note $A$ : \og le
nombre est un multiple de $4$\fg{} et $B$ : \og le nombre est un multiple de $6$\fg.
\begin{enumerate}
  \item Écrire $A$, $B$, $A\cap B$ et $A\cup B$ sous forme d'ensembles.
  \item Calculer $P(A)$, $P(B)$ et $P(A\cap B)$.
  \item Calculer $P(A\cup B)$ avec la formule, puis vérifier en comptant.
  \item Décrire $\overline{A}$ par une phrase et calculer $P(\overline{A})$.
\end{enumerate}

\exo[Avec les formules]
On choisit au hasard un élève d'un lycée. On note $S$ : \og il pratique un sport en
club\fg{} et $M$ : \og il joue d'un instrument de musique\fg. On sait que
$P(S)=0{,}45$, $P(M)=0{,}3$ et $P(S\cap M)=0{,}15$.
\begin{enumerate}
  \item Décrire par une phrase les événements $S\cap M$ et $S\cup M$.
  \item Calculer $P(S\cup M)$.
  \item Calculer la probabilité que l'élève ne fasse ni sport en club ni musique.
\end{enumerate}

\exo[Retrouver une probabilité]
Pour chaque cas, $A$ et $B$ sont deux événements d'une même expérience.
\begin{enumerate}[label=\alph*)]
  \item $P(A)=0{,}4$ ; $P(A\cup B)=0{,}7$ ; $P(A\cap B)=0{,}1$. Calculer $P(B)$, puis
    $P(\overline{B})$.
  \item $P(\overline{A})=0{,}35$ ; $P(\overline{B})=0{,}5$ ; $P(A\cup B)=0{,}8$. Calculer
    $P(A)$ et $P(B)$, puis $P(A\cap B)$.
  \item $A$ et $B$ sont incompatibles ; $P(A)=0{,}2$ et $P(B)=0{,}45$. Calculer
    $P(A\cup B)$.
\end{enumerate}

\exo[Deux activités dans une classe]
Dans une classe de $32$ élèves, $18$ pratiquent un sport, $10$ jouent d'un instrument et
$6$ font les deux. On choisit un élève au hasard.
\begin{enumerate}
  \item Représenter la situation par un diagramme (univers et deux événements), en
    écrivant l'effectif de chaque zone.
  \item Calculer la probabilité que l'élève pratique au moins l'une des deux activités.
  \item Calculer la probabilité qu'il n'en pratique aucune.
  \item Calculer la probabilité qu'il en pratique exactement une.
\end{enumerate}

\partiefiche{Probabilités conditionnelles et tableaux}

\exo[Calculer dans un tableau]
Une entreprise a interrogé ses $300$ salariés sur leur moyen de transport.
\begin{center}
\renewcommand{\arraystretch}{1.25}
\begin{tabular}{|l|c|c|c|}
\hline
 & Vélo & Autre & Total\\\hline
Moins de $30$ ans & $60$ & $90$ & $150$\\\hline
$30$ ans ou plus & $30$ & $120$ & $150$\\\hline
Total & $90$ & $210$ & $300$\\\hline
\end{tabular}
\end{center}
On choisit un salarié au hasard. On note $J$ : \og il a moins de $30$ ans\fg{} et
$V$ : \og il vient à vélo\fg.
\begin{enumerate}
  \item Calculer $P(J)$, $P(V)$ et $P(J\cap V)$.
  \item Le salarié a moins de $30$ ans. Calculer la probabilité qu'il vienne à vélo.
  \item Le salarié vient à vélo. Calculer la probabilité qu'il ait moins de $30$ ans.
\end{enumerate}

\exo[Compléter le tableau d'abord]
Sur $400$ candidats au permis de conduire, $220$ sont des femmes. $260$ candidats ont
réussi, dont $140$ femmes. On choisit un candidat au hasard ; on note $F$ : \og c'est
une femme\fg{} et $R$ : \og il a réussi\fg.
\begin{enumerate}
  \item Construire le tableau des effectifs ($F$, $\overline{F}$ ; $R$, $\overline{R}$).
  \item Calculer $P(\overline{F}\cap R)$, $P(F\cap\overline{R})$ et $P(\overline{R})$.
  \item Le candidat est un homme. Calculer la probabilité qu'il ait réussi.
  \item Le candidat a échoué. Calculer la probabilité que ce soit une femme.
\end{enumerate}

\exo[Trois colonnes]
On a demandé à $200$ élèves leur genre de série préféré.
\begin{center}
\renewcommand{\arraystretch}{1.25}
\begin{tabular}{|l|c|c|c|c|}
\hline
 & Comédie & Policier & Science-fiction & Total\\\hline
Seconde & $40$ & $30$ & $30$ & $100$\\\hline
Première & $20$ & $50$ & $30$ & $100$\\\hline
Total & $60$ & $80$ & $60$ & $200$\\\hline
\end{tabular}
\end{center}
On choisit un élève au hasard.
\begin{enumerate}
  \item Calculer la probabilité qu'il préfère les séries policières.
  \item L'élève est en première. Calculer la probabilité qu'il préfère les séries
    policières.
  \item L'élève préfère la science-fiction. Calculer la probabilité qu'il soit en seconde.
  \item Vrai ou faux : \og parmi les élèves qui préfèrent les comédies, deux sur trois
    sont en seconde\fg{} ? Justifier.
\end{enumerate}

\exo[Des pourcentages au tableau]
Un site de vente a reçu $2\,000$ commandes : $60\,\%$ passées sur le site internet, les
autres sur l'application. $5\,\%$ des commandes passées sur le site et $3\,\%$ de celles
passées sur l'application ont été renvoyées. On choisit une commande au hasard ; on
note $W$ : \og elle a été passée sur le site\fg{} et $R$ : \og elle a été
renvoyée\fg.
\begin{enumerate}
  \item Construire le tableau des effectifs.
  \item Calculer $P(R)$.
  \item La commande a été renvoyée. Calculer la probabilité qu'elle ait été passée sur
    le site (fraction irréductible).
\end{enumerate}

\partiefiche{Formule des probabilités conditionnelles}

\exo[Traduire, puis calculer]
On choisit au hasard un élève d'un lycée. On note $I$ : \og il est interne\fg{} et
$S$ : \og il est inscrit à l'association sportive\fg.
\begin{itemize}
  \item $30\,\%$ des élèves sont internes ;
  \item $12\,\%$ des élèves sont internes et inscrits à l'association sportive ;
  \item parmi les élèves qui ne sont pas internes, $25\,\%$ sont inscrits à
    l'association sportive.
\end{itemize}
\begin{enumerate}
  \item Traduire chaque donnée par une probabilité.
  \item Calculer $P_{I}(S)$ et interpréter le résultat.
  \item Calculer la probabilité que l'élève ne soit pas interne et soit inscrit à
    l'association sportive.
\end{enumerate}

\exo[Appliquer la formule]
Pour chaque cas, $A$ et $B$ sont deux événements.
\begin{enumerate}[label=\alph*)]
  \item $P(A)=0{,}5$ ; $P(B)=0{,}4$ ; $P(A\cap B)=0{,}1$. Calculer $P_{A}(B)$ et
    $P_{B}(A)$.
  \item $P(A)=0{,}6$ ; $P_{A}(B)=0{,}3$. Calculer $P(A\cap B)$ et $P_{A}(\overline{B})$.
  \item $P(B)=0{,}25$ ; $P(A\cap B)=0{,}05$. Calculer $P_{B}(A)$ et $P_{B}(\overline{A})$.
  \item $P(A)=\dfrac{1}{3}$ ; $P(A\cap B)=\dfrac{1}{12}$. Calculer $P_{A}(B)$.
  \item $P(A)=0{,}3$ ; $P_{A}(B)=0{,}8$ ; $P_{B}(A)=0{,}6$. Calculer $P(A\cap B)$,
    puis $P(B)$.
\end{enumerate}

\exo[\og Et\fg{} ou \og parmi\fg{} ?]
On choisit au hasard un inscrit d'une médiathèque. On note $F$ : \og il emprunte des
films\fg{} et $L$ : \og il emprunte des livres\fg.
\begin{enumerate}
  \item Traduire chaque phrase par une probabilité.
  \begin{enumerate}[label=\alph*)]
    \item $35\,\%$ des inscrits empruntent des films.
    \item Parmi les inscrits qui empruntent des films, $40\,\%$ empruntent aussi des
      livres.
    \item $14\,\%$ des inscrits empruntent des films et des livres.
    \item Un inscrit sur deux qui n'emprunte pas de films emprunte des livres.
  \end{enumerate}
  \item Vérifier que les données a) et b) donnent bien la donnée c).
  \item Calculer la probabilité qu'un inscrit n'emprunte pas de films mais emprunte des
    livres.
\end{enumerate}

\exo[Vrai ou faux]
$A$ et $B$ sont deux événements tels que $P(A)=0{,}4$, $P(B)=0{,}5$ et
$P(A\cap B)=0{,}2$. Pour chaque affirmation, dire si elle est vraie ou fausse, en
justifiant.
\begin{enumerate}
  \item $P_{A}(B)=0{,}5$.
  \item $P_{B}(A)=P_{A}(B)$.
  \item $P_{A}(\overline{B})=0{,}5$.
  \item $P(A\cup B)=0{,}9$.
\end{enumerate}

\partiefiche{Arbres pondérés}

\exo[Compléter un arbre]
On donne l'arbre pondéré ci-dessous.
\begin{center}
\begin{tikzpicture}[x=0.75cm,y=0.65cm]
  \arbvaleurs
  \arbre{A}{B}{0{,}4}{\dots}{0{,}7}{\dots}{\dots}{0{,}2}
\end{tikzpicture}
\end{center}
\begin{enumerate}
  \item Pour chaque nombre inscrit sur l'arbre, écrire la probabilité qu'il représente.
  \item Compléter l'arbre.
  \item Calculer $P(A\cap\overline{B})$ et $P(\overline{A}\cap B)$.
\end{enumerate}

\exo[Construire un arbre]
Sur une application de covoiturage, $70\,\%$ des trajets réservés sont des trajets
courts. Parmi les trajets courts, $20\,\%$ sont annulés ; parmi les autres, $10\,\%$. On
choisit un trajet au hasard ; on note $C$ : \og le trajet est court\fg{} et $N$ :
\og il est annulé\fg.
\begin{enumerate}
  \item Construire l'arbre pondéré qui décrit la situation.
  \item Calculer $P(C\cap N)$ et interpréter le résultat.
  \item Calculer $P(\overline{C}\cap\overline{N})$ et interpréter le résultat.
\end{enumerate}

\exo[Une branche par la formule]
$A$ et $B$ sont deux événements tels que $P(A)=0{,}25$, $P(A\cap B)=0{,}1$ et
$P_{\overline{A}}(B)=0{,}4$.
\begin{enumerate}
  \item Calculer $P_{A}(B)$.
  \item Construire l'arbre pondéré complet.
  \item Calculer $P(\overline{A}\cap\overline{B})$.
\end{enumerate}

\exo[Trouver l'erreur]
Dans un refuge, $60\,\%$ des animaux sont des chiens, les autres des chats. $18\,\%$ des
animaux sont des chiens adoptés dans le mois. Parmi les chats, $40\,\%$ sont adoptés dans
le mois. On note $C$ : \og l'animal est un chien\fg{} et $A$ : \og il est adopté dans le
mois\fg. Un élève a construit cet arbre :
\begin{center}
\begin{tikzpicture}[x=0.75cm,y=0.65cm]
  \arbvaleurs
  \arbre{C}{A}{0{,}6}{0{,}4}{0{,}18}{0{,}82}{0{,}4}{0{,}6}
\end{tikzpicture}
\end{center}
\begin{enumerate}
  \item Quelle erreur a-t-il faite ?
  \item Calculer $P_{C}(A)$, puis construire l'arbre correct.
  \item Calculer la probabilité qu'un animal soit un chat qui n'est pas adopté dans le
    mois.
\end{enumerate}

\partiefiche{Formule des probabilités totales}

\exo[Calculer $P(C)$, puis une probabilité inversée]
Dans une salle de sport, $40\,\%$ des adhérents viennent en semaine, les autres le
week-end. Parmi ceux qui viennent en semaine, $75\,\%$ prennent un cours collectif ;
parmi les autres, $45\,\%$. On choisit un adhérent au hasard ; on note $S$ : \og il
vient en semaine\fg{} et $C$ : \og il prend un cours collectif\fg.
\begin{enumerate}
  \item Construire l'arbre pondéré.
  \item Montrer que $P(C)=0{,}57$.
  \item L'adhérent prend un cours collectif. Calculer la probabilité qu'il vienne en
    semaine (fraction irréductible).
\end{enumerate}

\exo[Trois offres]
Un fournisseur d'énergie propose trois offres : $50\,\%$ de ses clients ont l'offre
$O_{1}$, $35\,\%$ l'offre $O_{2}$ et les autres l'offre $O_{3}$. Dans l'année, $4\,\%$ des
clients de l'offre $O_{1}$ résilient leur contrat, $6\,\%$ de ceux de $O_{2}$ et
$10\,\%$ de ceux de $O_{3}$. On choisit un client au hasard ; $R$ : \og il résilie\fg.
\begin{enumerate}
  \item Construire l'arbre pondéré.
  \item Calculer $P(R)$.
  \item Le client a résilié. Calculer la probabilité qu'il ait l'offre $O_{3}$ (fraction
    irréductible).
\end{enumerate}

\exo[Une probabilité inconnue]
Dans un lycée, $40\,\%$ des élèves sont demi-pensionnaires. Parmi eux, $60\,\%$ prennent
le bus. Au total, $51\,\%$ des élèves prennent le bus. On choisit un élève au hasard ;
$D$ : \og il est demi-pensionnaire\fg{} et $B$ : \og il prend le bus\fg. On note
$x=P_{\overline{D}}(B)$.
\begin{enumerate}
  \item Construire l'arbre pondéré, avec $x$.
  \item Montrer que $0{,}6x+0{,}24=0{,}51$.
  \item Résoudre cette équation et interpréter la solution.
\end{enumerate}

\exo[Un test de dépistage]
Une maladie touche $2\,\%$ d'une population. Un test est positif pour $95\,\%$ des
malades et pour $10\,\%$ des personnes qui ne sont pas malades. On choisit une personne
au hasard ; $M$ : \og elle est malade\fg{} et $T$ : \og son test est positif\fg.
\begin{enumerate}
  \item Construire l'arbre pondéré.
  \item Montrer que $P(T)=0{,}117$.
  \item Le test est positif. Calculer la probabilité que la personne soit malade
    (arrondie au centième). Ce résultat vous surprend-il ?
  \item Le test est négatif. Calculer la probabilité que la personne ne soit pas malade
    (arrondie au millième). Un test négatif est-il rassurant ?
\end{enumerate}

\partiefiche{Événements indépendants}

\exo[Indépendants ou non ?]
\begin{enumerate}
  \item On lance un dé équilibré à $8$ faces, numérotées de $1$ à $8$. On note $A$ :
    \og le nombre est pair\fg{} et $B$ : \og le nombre est inférieur ou égal à
    $4$\fg. $A$ et $B$ sont-ils indépendants ?
  \item Pour chaque arbre, $A$ et $B$ sont-ils indépendants ?
\end{enumerate}
\begin{center}
\begin{tikzpicture}[x=0.5cm,y=0.6cm]
  \arbvaleurs
  \arbre{A}{B}{0{,}3}{0{,}7}{0{,}6}{0{,}4}{0{,}6}{0{,}4}
  \node[font=\small] at (0,1.6) {a)};
\end{tikzpicture}\hfill
\begin{tikzpicture}[x=0.5cm,y=0.6cm]
  \arbvaleurs
  \arbre{A}{B}{0{,}5}{0{,}5}{0{,}2}{0{,}8}{0{,}4}{0{,}6}
  \node[font=\small] at (0,1.6) {b)};
\end{tikzpicture}
\end{center}

\exo[Retrouver $P(F)$]
Un appareil peut présenter un défaut d'aspect (événement $A$) et un défaut de
fonctionnement (événement $F$), de façon indépendante. On sait que $P(A)=0{,}1$ et que
la probabilité qu'il présente au moins l'un des deux défauts est $0{,}28$. On note
$p=P(F)$.
\begin{enumerate}
  \item Exprimer $P(A\cap F)$ en fonction de $p$.
  \item Montrer que $0{,}9p+0{,}1=0{,}28$.
  \item En déduire $P(F)$.
\end{enumerate}

\exo[Calculer avec l'indépendance]
Pour un week-end, la météo annonce de la pluie samedi avec une probabilité de $0{,}3$ et
dimanche avec une probabilité de $0{,}4$. On suppose ces deux événements, notés $A$ et
$B$, indépendants.
\begin{enumerate}
  \item Calculer la probabilité qu'il pleuve les deux jours.
  \item Calculer la probabilité qu'il pleuve au moins un des deux jours.
  \item Calculer la probabilité qu'il ne pleuve aucun des deux jours.
\end{enumerate}

\exo[Indépendance et second degré]
$A$ et $B$ sont deux événements tels que $P(A)=x$, $P(B)=1-x$ et
$P(A\cap B)=0{,}21$, où $x$ est un nombre de l'intervalle $\intff{0}{1}$.
\begin{enumerate}
  \item Exprimer $P(A)\times P(B)$ en fonction de $x$.
  \item Montrer que, si $A$ et $B$ sont indépendants, alors $x^{2}-x+0{,}21=0$.
  \item Résoudre cette équation, puis donner les valeurs possibles de $P(A)$ et de
    $P(B)$.
\end{enumerate}

\partiefiche{Trois façons de tirer deux objets}

\exo[Avec remise, sans remise, simultané]
Une urne contient $2$ boules rouges et $2$ boules vertes. On tire deux boules au hasard ;
$M$ : \og les deux boules ont la même couleur\fg. Calculer $P(M)$ quand :
\begin{enumerate}[label=\alph*)]
  \item on tire une boule, on la remet, puis on en tire une seconde ;
  \item on tire une boule, on la garde, puis on en tire une seconde ;
  \item on tire les deux boules en même temps.
\end{enumerate}

\exo[Deux bonbons]
Un sachet contient $3$ bonbons au citron et $2$ à la fraise. Léo prend un bonbon au
hasard et le mange, puis en prend un second. On note $C_{1}$ : \og le premier bonbon
est au citron\fg{} et $C_{2}$ : \og le second bonbon est au citron\fg.
\begin{enumerate}
  \item Construire l'arbre pondéré.
  \item Calculer la probabilité que les deux bonbons soient à la fraise.
  \item En déduire la probabilité qu'il mange au moins un bonbon au citron.
  \item Calculer la probabilité qu'il mange un bonbon de chaque parfum.
\end{enumerate}

\exo[Un code de deux chiffres]
On forme un code de deux chiffres pris au hasard parmi $1$, $2$, $3$, $4$. On note $E$ :
\og les deux chiffres sont pairs\fg.
\begin{enumerate}
  \item Un même chiffre peut être répété. Calculer $P(E)$.
  \item Les deux chiffres doivent être différents. Calculer $P(E)$.
  \item On choisit deux chiffres différents sans ordre (on liste les paires). Calculer
    $P(E)$.
  \item Pourquoi trouve-t-on le même résultat aux questions 2 et 3 ?
\end{enumerate}

\partiefiche{Épreuve de Bernoulli}

\exo[Trois tirs au but]
Une joueuse de football marque un tir au but avec une probabilité de $0{,}6$. Elle tire
trois tirs au but, de façon indépendante ; $S$ : \og le tir est marqué\fg.
\begin{enumerate}
  \item Construire l'arbre pondéré.
  \item Calculer la probabilité qu'elle marque les trois tirs.
  \item Calculer la probabilité qu'elle en marque exactement deux.
  \item Calculer la probabilité qu'elle en marque au moins un.
\end{enumerate}

\exo[Répondre au hasard]
Un QCM a trois questions ; chacune a quatre réponses, dont une seule est juste. Un élève
répond au hasard à chaque question.
\begin{enumerate}
  \item Expliquer pourquoi répondre à une question est une épreuve de Bernoulli, et
    donner son paramètre $p$.
  \item Construire l'arbre pondéré des trois réponses.
  \item Calculer la probabilité que les trois réponses soient justes.
  \item Calculer la probabilité qu'aucune réponse ne soit juste.
  \item Calculer la probabilité qu'au moins deux réponses soient justes.
\end{enumerate}

\exo[Deux feux tricolores]
Sur son trajet, Nora passe deux feux, qui fonctionnent indépendamment l'un de l'autre.
Chaque feu est vert à son arrivée avec la même probabilité $p$.
\begin{enumerate}
  \item Construire l'arbre pondéré, avec $p$.
  \item Exprimer en fonction de $p$ la probabilité que les deux feux soient verts.
  \item On sait que cette probabilité vaut $0{,}36$. Déterminer $p$.
  \item Calculer la probabilité qu'un seul des deux feux soit vert.
\end{enumerate}

\partiefiche{Variable aléatoire}

\exo[Un jeu de dé]
On lance un dé équilibré à six faces. On gagne $4$\,€ si on obtient $5$ ou $6$, on gagne
$1$\,€ si on obtient $4$, et on perd $3$\,€ sinon. On note $X$ le gain, en euros.
\begin{enumerate}
  \item Donner les valeurs prises par $X$.
  \item Calculer $P(X=4)$ et $P(X=-3)$.
  \item Calculer $P(X\ge 1)$.
  \item Décrire par une phrase l'événement $\{X\le 1\}$ et calculer sa probabilité.
\end{enumerate}

\exo[Un cube peint]
\begin{minipage}[t]{0.68\linewidth}\raggedright
Un grand cube en bois, peint sur ses six faces, est découpé en $27$ petits cubes
identiques. On tire un petit cube au hasard ; $X$ est le nombre de ses faces peintes.
\end{minipage}\hfill
\begin{minipage}[t]{0.28\linewidth}\vspace{-2pt}
\begin{tikzpicture}[scale=0.32]
  \fill[coulExo!25] (0,0) rectangle (3,3);
  \fill[coulExo!15] (0,3) -- (3,3) -- (4.2,3.9) -- (1.2,3.9) -- cycle;
  \fill[coulExo!35] (3,0) -- (4.2,0.9) -- (4.2,3.9) -- (3,3) -- cycle;
  \draw (0,0) grid (3,3);
  \foreach \i in {0,1,2,3}{
    \draw (\i,3) -- ++(1.2,0.9);
    \draw (0.4*\i,3+0.3*\i) -- ++(3,0);
    \draw (3,\i) -- ++(1.2,0.9);
    \draw (3+0.4*\i,0.3*\i) -- ++(0,3);}
\end{tikzpicture}
\end{minipage}
\begin{enumerate}
  \item Combien de petits cubes ont trois faces peintes ? deux ? une seule ? aucune ?
  \item Donner les valeurs prises par $X$.
  \item Calculer $P(X=2)$ et $P(X\ge 1)$.
  \item Décrire par une phrase l'événement $\{X\ge 2\}$ et calculer sa probabilité.
\end{enumerate}

\exo[La somme de deux jetons]
Une urne contient quatre jetons numérotés $1$, $2$, $3$, $4$. On tire deux jetons en
même temps ; $X$ est la somme des deux numéros.
\begin{enumerate}
  \item Écrire les $6$ paires possibles et la valeur de $X$ pour chacune.
  \item Donner les valeurs prises par $X$.
  \item Calculer $P(X=5)$ et $P(X\ge 6)$.
  \item Calculer la probabilité que $X$ soit un nombre pair.
\end{enumerate}

\partiefiche{Loi de probabilité}

\exo[Un jeu avec mise]\label{exo:mise}
Pour jouer, on paie $3$\,€. Un sac contient $2$ jetons gagnants sur $10$. On tire un
jeton, on le remet, puis on en tire un second. Chaque jeton gagnant rapporte $5$\,€. On
note $X$ le gain algébrique : ce que l'on reçoit, moins les $3$\,€ payés.
\begin{enumerate}
  \item Donner les valeurs prises par $X$.
  \item Déterminer la loi de probabilité de $X$.
\end{enumerate}

\exo[Compléter une loi]
\begin{enumerate}
  \item La loi de $Y$ est donnée ci-dessous. Calculer $a$, puis $P(Y\ge 0)$ et $P(Y>0)$.
\begin{center}
\renewcommand{\arraystretch}{1.3}
\begin{tabular}{|c|c|c|c|c|}
\hline
$y_{i}$ & $-2$ & $0$ & $1$ & $4$\\\hline
$P(Y=y_{i})$ & $0{,}15$ & $a$ & $0{,}4$ & $0{,}1$\\\hline
\end{tabular}
\end{center}
  \item La loi de $Z$ est donnée ci-dessous. Calculer $a$.
\begin{center}
\renewcommand{\arraystretch}{1.3}
\begin{tabular}{|c|c|c|c|}
\hline
$z_{i}$ & $1$ & $2$ & $3$\\\hline
$P(Z=z_{i})$ & $0{,}4$ & $a$ & $2a$\\\hline
\end{tabular}
\end{center}
\end{enumerate}

\exo[Le prix d'une entrée]
Dans un parc d'attractions, $70\,\%$ des visiteurs ont un pass annuel. Parmi eux,
$30\,\%$ sont des enfants ; parmi les visiteurs sans pass, $50\,\%$. Un visiteur sans pass
paie $30$\,€ ; un adulte avec pass paie $10$\,€ ; un enfant avec pass ne paie pas. On
choisit un visiteur au hasard ; $X$ est le prix qu'il paie, en euros.
\begin{enumerate}
  \item Construire l'arbre pondéré (événements $A$ : \og il a un pass\fg{} et $E$ :
    \og c'est un enfant\fg).
  \item Donner les valeurs prises par $X$.
  \item Déterminer la loi de probabilité de $X$.
\end{enumerate}

\exo[Un dé truqué]
Un dé à six faces est truqué : la probabilité d'obtenir la face $k$ est égale à
$k\times q$, où $q$ est un nombre positif.
\begin{enumerate}
  \item Montrer que $21q=1$.
  \item Déterminer la loi de probabilité du numéro obtenu.
  \item Calculer la probabilité d'obtenir un nombre pair.
\end{enumerate}

\partiefiche{Espérance, variance, écart-type}

\exo[Le jeu est-il favorable ?]\label{exo:favorable}
On reprend le jeu de l'exercice~\ref{exo:mise} (mise de $3$\,€) :
\begin{center}
\renewcommand{\arraystretch}{1.3}
\begin{tabular}{|c|c|c|c|}
\hline
$x_{i}$ & $-3$ & $2$ & $7$\\\hline
$P(X=x_{i})$ & $0{,}64$ & $0{,}32$ & $0{,}04$\\\hline
\end{tabular}
\end{center}
\begin{enumerate}
  \item Calculer $E(X)$.
  \item Interpréter le résultat. Le jeu est-il favorable au joueur ?
  \item L'organisateur remplace la mise de $3$\,€ par une mise de $m$ euros. Quelle mise
    rend le jeu équitable ?
\end{enumerate}

\exo[Variance et écart-type]
Pour le jeu de l'exercice~\ref{exo:favorable} (mise de $3$\,€), on a trouvé $E(X)=-1$. Calculer $V(X)$,
puis $\sigma(X)$ arrondi au centième.

\exo[Deux jeux à comparer]
Jeu A : on gagne $2$\,€ avec une probabilité de $0{,}5$, on perd $2$\,€ sinon ; gain
$X_{A}$.\par
Jeu B : on gagne $18$\,€ avec une probabilité de $0{,}1$, on perd $2$\,€ sinon ; gain
$X_{B}$.
\begin{enumerate}
  \item Calculer $E(X_{A})$ et $E(X_{B})$.
  \item Calculer $\sigma(X_{A})$ et $\sigma(X_{B})$.
  \item Lequel des deux jeux est le plus risqué ? Justifier.
\end{enumerate}

\exo[La tombola de la fête du lycée]
Pour la fête du lycée, on vend $200$ billets de tombola à $2$\,€. Un billet gagne un lot
de $100$\,€, quatre billets un lot de $20$\,€ et dix billets un lot de $5$\,€. On achète un
billet au hasard ; $X$ est le gain algébrique (lot moins prix du billet).
\begin{enumerate}
  \item Donner les valeurs prises par $X$.
  \item Déterminer la loi de probabilité de $X$.
  \item Calculer $E(X)$ et interpréter le résultat.
  \item Avec les mêmes lots et le même prix, combien de billets faudrait-il vendre pour
    que le jeu soit équitable ?
\end{enumerate}

\columnbreak
\exo[Une valeur inconnue]
Une variable aléatoire $X$ prend les valeurs $-2$ ; $a$ ; $3$ avec les
probabilités $0{,}3$ ; $0{,}5$ ; $0{,}2$, où $a$ est un nombre.
\begin{enumerate}
  \item Exprimer $E(X)$ en fonction de $a$.
  \item On sait que $E(X)=2$. Calculer $a$.
  \item Calculer $V(X)$, puis $\sigma(X)$ arrondi au centième.
\end{enumerate}

\partiefichesansnum{Exercices type bac}

\exo[Automatismes (QCM)]
Pour chaque question, une seule réponse est exacte.
\begin{enumerate}
  \item On lance un dé équilibré à six faces. La probabilité d'obtenir un nombre
    strictement supérieur à $4$ est :\par
    a)~$\frac{1}{6}$ \quad b)~$\frac{1}{3}$ \quad c)~$\frac{1}{2}$ \quad d)~$\frac{2}{3}$
  \item Une classe compte $12$ filles, dont $3$ internes. On choisit une fille au
    hasard. La probabilité qu'elle soit interne est :\par
    a)~$\frac{1}{10}$ \quad b)~$\frac{1}{3}$ \quad c)~$\frac{1}{4}$ \quad d)~$\frac{2}{5}$
  \item $P(A)=0{,}3$ ; $P_{A}(B)=0{,}2$ ; $P_{\overline{A}}(B)=0{,}5$. Alors
    $P(\overline{A}\cap B)$ vaut :\par
    a)~$0{,}06$ \quad b)~$0{,}35$ \quad c)~$0{,}5$ \quad d)~$0{,}41$
  \item $A$ et $B$ sont indépendants, $P(A)=0{,}5$ et $P(B)=0{,}4$. Alors
    $P(A\cup B)$ vaut :\par
    a)~$0{,}9$ \quad b)~$0{,}2$ \quad c)~$0{,}7$ \quad d)~$0{,}1$
  \item $X$ prend les valeurs $-1$ ; $2$ ; $5$ avec les probabilités $0{,}5$ ; $0{,}3$ ;
    $0{,}2$. Alors $E(X)$ vaut :\par
    a)~$1{,}1$ \quad b)~$2$ \quad c)~$0{,}3$ \quad d)~$1{,}5$
\end{enumerate}

\exo[Location de kayaks]
Une base nautique loue des kayaks. On sait que :
\begin{itemize}
  \item $70\,\%$ des locations ont lieu en semaine, les autres le week-end ;
  \item parmi les locations en semaine, $20\,\%$ durent plus d'une heure ;
  \item $23\,\%$ de l'ensemble des locations durent plus d'une heure.
\end{itemize}
On choisit une location au hasard ; $S$ : \og elle a lieu en semaine\fg{} et $L$ :
\og elle dure plus d'une heure\fg.
\begin{enumerate}
  \item Recopier et compléter les pointillés de l'arbre.
\begin{center}
\begin{tikzpicture}[x=0.75cm,y=0.65cm]
  \arbvaleurs
  \arbre{S}{L}{\dots}{\dots}{\dots}{\dots}{}{}
\end{tikzpicture}
\end{center}
  \item Donner, par simple lecture de l'énoncé, $P(L)$.
  \item Montrer que $P(S\cap L)=0{,}14$, puis en déduire que $P(\overline{S}\cap L)=0{,}09$.
  \item Calculer la probabilité qu'une location du week-end dure plus d'une heure.
  \item Une location a duré plus d'une heure. Calculer la probabilité qu'elle ait eu lieu
    le week-end (fraction irréductible).
  \item Les événements $S$ et $L$ sont-ils indépendants ? Justifier.
\end{enumerate}

\exo[Un jeu à la fête foraine]
Une urne contient $4$ boules rouges et $6$ vertes. On paie $2$\,€ pour tirer une boule,
la remettre et en tirer une seconde ; chaque boule rouge tirée rapporte $3$\,€. On note
$R_{1}$ (et $R_{2}$) : \og la première (la seconde) boule est rouge\fg, et $X$ le gain
algébrique.
\begin{enumerate}
  \item Construire l'arbre pondéré.
  \item Donner les valeurs prises par $X$.
  \item Montrer que $P(X=1)=0{,}48$ ; donner la loi de $X$.
  \item Calculer $E(X)$ et interpréter le résultat.
  \item Quelle mise rendrait le jeu équitable ?
\end{enumerate}

\end{multicols}

\end{document}
